\documentclass{subfiles}
\begin{document}
    Now that we have seen a couple of public key encryption algorithm,
        a question should pop inside the reader's head:
        what garantees do we have about these algorithms?
    
    It should come to no surprise that the answer is given by math.
        Precisely, by the hardness of the problems underlying these algorithms.
        For instance, the security of RSA is given by the hard problem of integers factorization,
        which asks us, given an integer \(n\), to find its factorization into primes.
        The Diffie-Hellman algorithm, and therefore ElGamal,
        are based on the \emph{discrete algorithm} problem,
        which roughuly asks, given and integer \(\alpha\), 
        to find \(a\) such that \(\alpha \equiv \gamma^{a} \mod\), where \(\gamma \text{ and } p\) are known.

    These we presented are just a couple of several public key encryption algorithm.
        The interested reader can easely find a whole literature out there.
\end{document}
